\documentclass[a4paper,11pt,dvipsnames]{article}
\usepackage{amsmath,amsthm,amssymb,amsfonts}
\usepackage{array}
\usepackage{tabularx}
\usepackage[table]{xcolor}
\usepackage[most]{tcolorbox}
\usepackage{graphicx}
\graphicspath{{./img/}}
\usepackage[a4paper,left=1.5cm,right=1.5cm,top=1cm,bottom=1.5cm]{geometry}
% \usepackage[hmargin=1cm,vmargin=1cm]{geometry}
\usepackage{array}
\usepackage{setspace}
\newcommand*\diff{\mathop{}\!\mathrm{d}}
\newcommand*\Diff[1]{\mathop{}\!\mathrm{d^#1}}

\newcommand{\N}{\mathbb{N}}
\newcommand{\Z}{\mathbb{Z}}

% A style for tcolorbox without frame etc. "notcb"
\tcbset{notcb/.style={enhanced jigsaw, colback=white, coltitle=black,sharp corners, boxrule=0pt}}

% Style for the number tables
\tcbset{tabularheader/.style={arc=2pt,notcb,boxsep=0pt,left=0pt,colbacktitle=white,rounded corners=all,enhanced jigsaw,coltitle=black,tabularx={*{6}{Y|}Y},boxrule=1pt,attach boxed title to top left, boxed title style={enhanced jigsaw,sharp corners,left=0pt,boxrule=0pt}}}

\newcolumntype{Y}{>{\centering\arraybackslash}X}

\usepackage[locale=FR, per-mode=fraction, separate-uncertainty=true]{siunitx}
\usepackage{physics}
\usepackage[utf8]{inputenc}
\usepackage[T1]{fontenc}
\usepackage[spanish]{babel}

\usepackage{enumitem}

\usepackage{pgfplots}
% \usepackage{pdfpages}
\usepackage{adjustbox, tikz}
\usepackage{circuitikz}
\usepackage{tikz-3dplot}

% \usepackage{gnuplottex}
\usepackage{xcolor}
\usepackage[locale=FR, per-mode=fraction, separate-uncertainty=true]{siunitx}
\usepackage{physics}

\usepackage{ifpdf}
\ifpdf%
  \usepackage{pdftricks}
  \begin{psinputs}
    \usepackage{pst-optic}
%\usepackage{pstricks-add} 	
  \end{psinputs}
\else
  \usepackage{pst-optic}
  \newenvironment{pdfpic}{}{}
\fi
 

\usetikzlibrary{decorations}
\usetikzlibrary{decorations.pathmorphing}
\usetikzlibrary{shapes.geometric}
\usetikzlibrary{arrows}
\usetikzlibrary{arrows.meta}   %edg
\usetikzlibrary{calc}
\pgfplotsset{compat=newest}
\tikzset{axis/.style={blue, thick,-latex}}


\begin{document}

{\Large Resoluciones 3er parcial - 2022}

\begin{enumerate}[leftmargin=*]

\item 
% \begin{minipage}[t]{.6\textwidth}
  Considerando que el capacitor en el circuito mostrado en la figura se encuentra descargado, que las resistencias son idénticas y que en el instante inicial se cierra la llave, seleccione la opción que mejor describe el valor de la corriente $i$ una vez que se alcanza el régimen estacionario, comparado con su valor inicial (inmediatamente después de cerrar la llave).

\begin{minipage}[t]{.65\textwidth}
\begin{tabularx}{\textwidth}{XX}
$\square$ 25\% menor que la corriente inicial. & $\square$ 25\% mayor que la corriente inicial.\\
$\square$ $33.3$\% menor que la corriente inicial. & $\square$ $33.3$\% mayor que la corriente inicial.\\
$\square$ $50$\% menor que la corriente inicial. & $\square$ $50$\% mayor que la corriente inicial.\\
$\square$ No circula corriente en el estado estacionario.&\\
\end{tabularx}  
\end{minipage}
\begin{minipage}[t]{.3\textwidth}
% \strut\vspace*{-\baselineskip}
\strut\vspace*{-90pt}
\begin{center}
  \begin{circuitikz}[scale=1]
    \draw [] (0,0) to[R=$R$] (0,-1.8) to[R=$R$] (2,-1.8);
    \draw [] (0,0) to[battery2, l=$\varepsilon$, invert] (2,0) to[switch]  (4.5,0) -- (4.5,-1.8) -- (4,-1.8);
    \draw [] (2,-1.3) to[R=$R$] (4,-1.3) -- (4,-2.3) to[C, l=$C$] (2,-2.3) -- (2,-1.3);
    \draw [-Stealth] (-0.4, -1.5) -- (-0.4, -0.2) node [midway, left] {$i$};
  \end{circuitikz}
\end{center}
\end{minipage}

\

\textit{Resolución:}\\

En el instante inicial no circula corriente a través de la $R$ en paralelo con el capacitor, entonces la corriente inicial vale:
\begin{flalign*}
    i &= \frac{\varepsilon}{2R}
\end{flalign*}

En el estado estacionario, no hay corrientes en la rama del capacitor, solo circula a través de la $R$ en paralelo. La corriente en el estado estacionario es:
\begin{flalign*}
    i' &= \frac{\varepsilon}{3R}
\end{flalign*}

Por lo tanto: $i' = \frac{2}{3}i$

Respuesta: $33.3$\% menor que la corriente inicial.

\vspace*{1cm}

\begin{minipage}[t]{.6\textwidth}
\item En la figura se muestran dos cables infinitos paralelos, perpendiculares al plano $xy$, uno transportando una corriente $i_1$ y el otro transportando una corriente $i_2$. El campo magnético producido por esas corrientes en el punto $P$ es:
\begin{align*}
  \vec{B} = \left ( -23.3\,\hat{x} + 8.5\,\hat{y} \right ) \, \si{\micro\tesla}
\end{align*}
Hallar la intensidad y el sentido de las corrientes $i_1$ e $i_2$.
\end{minipage}
\begin{minipage}[t]{.4\textwidth}
\strut\vspace*{-\baselineskip}
\begin{center}
  \begin{tikzpicture}[scale=0.4]
    \draw [-{Stealth}] (0,-1) -- (0,8) node [left] {$\hat{y}$};
    \draw [-{Stealth}] (-1,0) -- (5.5,0) node [below] {$\hat{x}$};
    % \draw [] (5+0.2828,-0.2828) -- (5-0.2828,0.2828);
    \fill [] (4,0) circle (0.3);
    % \fill [] (0,0) circle (0.2);
    \filldraw [fill=white] (0,0) circle (0.3);
    \filldraw [fill=white] (0,6) circle (0.3);
    % \draw [dotted] (0,0) -- (5,0) -- (0,3.33) -- (0,0);
    \draw [] (.8,-0.5) node [] {$i_1$};
    \draw [] (.8,6) node [] {$i_2$};
    \draw [] (4,0.7) node [] {$P$};

    \draw [dotted] (-2,6) -- (-0.3,6);
    \draw [dotted] (-2,0) -- (-0.3,0);
    \draw [|-|] (-2,0) -- (-2,6) node [midway, left] {$\SI{3.0}{\centi\metre}$};
    \draw [dotted] (0,-0.3) -- (0,-1.5);
    \draw [dotted] (4,-0.3) -- (4,-1.5);
    \draw [|-|] (0,-1.5) -- (4,-1.5) node [midway, below] {$\SI{2.0}{\centi\metre}$};
  \end{tikzpicture}
\end{center}
\end{minipage}


\

\textit{Resolución:}\\

En la siguiente figura se muestran las direcciones de los campos producidos por cada corriente. Se puede ver que el campo producido por $i_1$ no tiene componente en $x$, y solo el campo producido por $i_2$ tiene componente en $x$.

\begin{center}
  \begin{tikzpicture}[scale=0.4]
    \draw [-{Stealth}] (0,-1) -- (0,8) node [left] {$\hat{y}$};
    \draw [-{Stealth}] (-1,0) -- (5.5,0) node [below] {$\hat{x}$};
    % \draw [] (5+0.2828,-0.2828) -- (5-0.2828,0.2828);
    \fill [] (4,0) circle (0.3);
    % \fill [] (0,0) circle (0.2);
    \filldraw [fill=white] (0,0) circle (0.3);
    \filldraw [fill=white] (0,6) circle (0.3);
    % \draw [dotted] (0,0) -- (5,0) -- (0,3.33) -- (0,0);
    \draw [] (.8,-0.5) node [] {$i_1$};
    \draw [] (.8,6) node [] {$i_2$};
    \draw [] (4,0.7) node [] {$P$};

    \draw [dotted] (-2,6) -- (-0.3,6);
    \draw [dotted] (-2,0) -- (-0.3,0);
    \draw [dotted] (0,-0.3) -- (0,-1.5);
    \draw [dotted] (4,-0.3) -- (4,-1.5);
    \draw [red] (4,-3) -- (4,3) node [red, right] {1};
    \draw [blue] (1,-2) -- (7,2) node [blue, right] {2};
  \end{tikzpicture}
\end{center}

Entonces se puede afirmar que la componente $x$ del campo $B_2$ es igual a la componente $x$ del campo resultante en $P$, y como debe ser negativa, el sentido del vector del campo producido por $i_2$ en el punto $P$ es el siguiente:
\begin{center}
  \begin{tikzpicture}[scale=0.4]
    \draw [-{Stealth}] (0,-1) -- (0,8) node [left] {$\hat{y}$};
    \draw [-{Stealth}] (-1,0) -- (5.5,0) node [below] {$\hat{x}$};
    % \draw [] (5+0.2828,-0.2828) -- (5-0.2828,0.2828);
    \fill [] (4,0) circle (0.3);
    % \fill [] (0,0) circle (0.2);
    \filldraw [fill=white] (0,0) circle (0.3);
    \filldraw [fill=white] (0,6) circle (0.3);
    % \draw [dotted] (0,0) -- (5,0) -- (0,3.33) -- (0,0);
    \draw [] (.8,-0.5) node [] {$i_1$};
    \draw [] (.8,6) node [] {$i_2$};
    \draw [] (4,0.7) node [] {$P$};

    \draw [dotted] (-2,6) -- (-0.3,6);
    \draw [dotted] (-2,0) -- (-0.3,0);
    \draw [dotted] (0,-0.3) -- (0,-1.5);
    \draw [dotted] (4,-0.3) -- (4,-1.5);
    % \draw [red] (4,-3) -- (4,3) node [red, right] {1};
    \draw [blue, {Stealth}-] (1,-2) -- (4,0) node [blue, pos=0,left] {$\vec{B}_2$};
    \draw [blue] (2.5,-0.4) node [blue] {$\theta$};
    \draw [red, {Stealth}-{Stealth}] (0,6) -- (4,0) node [red, midway, right] {$d=\SI{3.61}{\centi\metre}$};
  \end{tikzpicture}
\end{center}

El ángulo mostrado en la figura es:
\begin{flalign*}
    \tan\theta &= \frac{2}{3}\\
    \theta &= \SI{33.7}{\degree}
\end{flalign*}

Igualando la componente $x$ de $B_2$ a la componente $x$ del campo resultante se tiene:
\begin{flalign*}
    B_{2,x} &= \SI{-23.3}{\micro\tesla}\\
    -|\vec{B}_2| \cos\SI{33.7}{\degree} &= \SI{-23.3}{\micro\tesla}\\
    |\vec{B}_2| &= \SI{2.8E-5}{\tesla}
\end{flalign*}

Ahora que se conoce el módulo del campo $B_2$ se puede calcular la corriente $i_2$:
\begin{flalign*}
    \frac{\mu_0 i_2}{2\pi d} &= \SI{2.8E-5}{\tesla}\\
    i_2 &= \frac{\SI{2.8E-5}{\tesla} \cdot \SI{0.0361}{\metre}}{\SI{2E-7}{\tesla.\metre/\ampere}}\\
    i_2 &= \SI{5.05}{\ampere}
\end{flalign*}

A continuación se busca el módulo del campo $B_1$ producido por $i_1$, utilizando la componente $y$ del campo resultante:
\begin{flalign*}
    B_y &= B_{1,y} + B_{2,y}\\
    \SI{8.5E-6}{\tesla} &= B_{1,y} - |\vec{B}_2|\sin\theta\\
    \SI{8.5E-6}{\tesla} &= B_{1,y} - \SI{2.8E-5}{\tesla}\cdot\SI{33.7}{\degree}\\
    B_{1,y} &= \SI{2.4E-5}{\tesla}
\end{flalign*}

El campo $B_1$ está dirigido hacia arriba, y su módulo es igual a la componente $y$ recién encontrada:
\begin{flalign*}
    |\vec{B}_1| &= B_{1,y}\\
    \frac{\mu_0 i_1}{2\pi \cdot \SI{0.02}{\metre}} &= \SI{2.4E-5}{\tesla}\\
    i_1 &= \SI{2.4}{\ampere}
\end{flalign*}

Respuestas: $i_1 = \SI{2.4}{\ampere}$ saliendo de la página; $i_2 = \SI{5.1}{\ampere}$ entrando en la página. 

\vspace*{1cm}

\begin{minipage}[t]{.5\textwidth}
\item
  Una espira cuadrada de 120 vueltas se encuentra inmersa en un campo magnético uniforme $\vec{B} = \SI{0.02}{\tesla}\,\hat{j}$. La espira tiene $\SI{3.0}{\centi\metre}$ de lado y forma un ángulo $\alpha$ variable con el plano $ik$. \textit{a}) Si se hace girar la espira alrededor del eje $k$ con una velocidad angular de $120\pi\si{s^{-1}}$, ¿cuánto vale la fem inducida en la espira en un instante en que $\alpha$ es $\pi/2$? \textit{b}) ¿En cuáles valores de $\alpha$ la corriente inducida invierte su sentido?
\end{minipage}
\begin{minipage}[t]{0.5\textwidth}
\strut\vspace*{-\baselineskip}
% \strut\vspace*{2.5cm}
\begin{center}
    \tdplotsetmaincoords{70}{120}
    \begin{tikzpicture}[tdplot_main_coords, scale=0.7]
        \draw[-{Stealth}] (0,0,0) -- (4,0,0) node [pos=1.1] {$i$};
        \draw[-{Stealth}] (0,0,0) -- (0,4,0) node [pos=1.05] {$j$};
        \draw[-{Stealth}] (0,0,0) -- (0,0,4)  node [left] {$k$};  
        \draw[] (0,0,0) -- (2,3,0) -- (2,3,3) -- (0,0,3) -- cycle;
        \draw[] (0.1,-0.05,0) -- (2.1,2.95,0) -- (2.1,2.95,3) -- (0.1,-0.05,3) -- cycle;
        \draw[] (-0.1,0.05,0) -- (1.9,3.05,0) -- (1.9,3.05,3) -- (-0.1,0.05,3) -- cycle;
        \draw[-{latex}, thick] (0.9,-2.5,1.25) -- (0.9,-0.5,1.25);
        \draw[-{latex}, thick] (0.9,3,1.25) -- (0.9,5,1.25) node [midway, above] {$\vec{B}$};
        \draw[-{latex}, thick] (0.9,-2.5,2.5) -- (0.9,-0.5,2.5);
        \draw[-{latex}, thick] (0.9,3,2.5) -- (0.9,5,2.5);
        \draw[-{latex}, thick] (0.9,-2.5,0) -- (0.9,-0.5,0);
        \draw[-{latex}, thick] (0.9,3,0) -- (0.9,5,0);
        \draw[-latex] (0:2.5) arc (0:53:2.5) node[black,midway,below] {$\alpha$};
    \end{tikzpicture}
\end{center}
\end{minipage}

\

\textit{Resolución:}\\
\textit{a}) Vista desde arriba: \\
\begin{center}
  \begin{tikzpicture}[scale=0.4]
    \draw [-{Stealth}] (0,1) -- (0,-5) node [left] {$\hat{i}$};
    \draw [-{Stealth}] (-1,0) -- (5,0) node [above] {$\hat{j}$};
    \draw [thick] (0,0) -- (6,-6);
    \draw [] (0.5,-1.5) node [] {$\alpha$};
    \draw [-{Stealth}] (3,-3) -- (7,1) node [right] {$\vec{A}$};
    \draw [-{Stealth}] (3,-3) -- (7,-3) node [below] {$\vec{B}$};
    \draw [] (4.5,-2.5) node [] {$\alpha$};
\end{tikzpicture}
\end{center}
\begin{flalign*}
    \Phi &= NAB\cos\alpha\\
\end{flalign*}
\begin{flalign*}
    \varepsilon &= -\frac{d\Phi}{dt}\\
    \varepsilon &= NAB\sin\alpha \frac{d\alpha}{dt}\\
\end{flalign*}
La velocidad angular es: $d\alpha/dt = 120\pi\,\si{s^{-1}}$. Entonces, la fem inducida cuando $\alpha = \pi/2$ es:
\begin{flalign*}
    \varepsilon &= 120\cdot (\SI{0.03}{\metre})^2 \cdot \SI{0.02}{\tesla} \cdot \sin(\pi/2) \cdot 120\pi\,\si{s^{-1}}\\
    \varepsilon &= \SI{0.814}{\volt}
\end{flalign*}

\textit{b}) La corriente inducida es $\varepsilon/R = \frac{NAB\omega}{R}\sin\alpha$.\\
Entonces cambia de signo cada vez que la función seno cambie de signo.\\

Respuesta: $\alpha = m\pi$ con $m \in \mathbb{Z}$.

\vspace*{1cm}

\item 
\begin{minipage}[t]{.6\textwidth}
  En el circuito de la figura, la diferencia de potencial en la batería es $\varepsilon = \SI{12.0}{\volt}$, la capacidad del capacitor es $C = \SI{2.4}{\milli\farad}$ y la resistencia $R_1$ vale $\SI{250}{\ohm}$. Con el capacitor inicialmente descargado, se hace contacto con la llave en la posición 1 y se espera un segundo, para luego mover la llave a la posición 2. Tres segundos después, con la llave aun en la posición 2, se mide la diferencia de potencial en las placas del capacitor y su valor en ese instante resulta $\SI{1.45}{\volt}$. ¿Cuánto vale la resistencia $R_2$?  
\end{minipage}
\begin{minipage}[t]{.4\textwidth}
\strut\vspace*{-\baselineskip}
\begin{center}
\begin{circuitikz}[scale=1]
    \draw (0.5,3.5) node[spdt, rotate=90] (Sw) {} (Sw.in) node[left] {} (Sw.out 1) node[below left] {1} (Sw.out 2) node[below right] {2};
    \draw [] (0.5,1.5) to[C, l=$C$] (Sw.in);
    % \draw [] (0,0) to[R=$R_2$, color=red] (0,2) to[C, l=$C$] (Sw.in) ;
    \draw [] (Sw.out 1) -| (-2,4) to[battery2, l=$\varepsilon$] (-2,1.5)  to[R=$R_1$] (0,1.5) -- (2,1.5) to[R=$R_2$] (2,4) |- (Sw.out 2);
\end{circuitikz}
\end{center}
\end{minipage}

\

\textit{Resolución:}\\

Con la llave en la posición 1: $\tau = \SI{250}{\ohm}\cdot \SI{2.4E-3}{\farad} = \SI{0.6}{\second}$.\\
La diferencia de potencial en el capacitor 1 segundo después es:
\begin{flalign*}
    V(t) &= V_{max} \left (1 - \text{e}^{-t/\tau} \right )\\
    V(\SI{1}{\second}) &= \SI{12}{\volt} \left (1 - \text{e}^{-1/0.6} \right ) = \SI{9.37}{\volt}
\end{flalign*}

Con la llave en la posición 2: $\tau = R_2C$.\\
La diferencia de potencial 3 segundos después de conectar en la posición 2:
\begin{flalign*}
    V(t) &= V_0 \text{e}^{-t/\tau}\\
    V(\SI{3}{\second}) = \SI{1.45}{\volt} &= \SI{9.37}{\volt} \text{e}^{-3/\tau}\\
    \tau &= \SI{1.58}{\second}
\end{flalign*}
Con ese valor del tiempo característico en la segunda etapa se puede encontrar la resistencia desconocida:
\begin{flalign*}
    \tau &= R_2C\\ 
    \SI{1.58}{\second} &= R_2 \cdot \SI{2.4E-3}{\farad}
\end{flalign*}

Respuesta: $R_2 = \SI{658}{\ohm}$.

\end{enumerate}

\end{document}
